\chapter{Display LCD}

Il display 16x2 usa un layer di astrazione definito in \texttt{display.h} (classe \texttt{DisplayInterface}) e la logica applicativa in \texttt{lcd.h} (classe \texttt{DisplayPanel}). Questa separazione permette di testare la logica senza dipendere direttamente dall'hardware.

\section{Pagine}
Il pannello ruota tre pagine ogni 5 secondi:
\begin{itemize}
  \item \textbf{Data/Ora}: calcolata a partire da \texttt{\_\_DATE\_\_}, \texttt{\_\_TIME\_\_} e \texttt{millis()}.
  \item \textbf{Stato}: temperatura, umidita e conteggio persone.
  \item \textbf{Diagnostica}: messaggi di errore o anomalie.
\end{itemize}
La pagina errori mostra in modo prioritario le condizioni critiche, facilitando l'intervento rapido.

\section{Calcolo data e ora}
Non essendo presente un RTC, l'orario e calcolato combinando il timestamp di compilazione con l'incremento di \texttt{millis()}. La base e determinata al costruttore, cosi da offrire una data iniziale coerente e un avanzamento temporale continuo.

\section{Aggiornamento}
Il refresh e limitato a 1 secondo per evitare flicker. Ogni update legge i valori dai moduli \texttt{Thermostat}, \texttt{HumidifierControl} e \texttt{LightingControl}. La pagina cambia ogni 5 secondi per lasciare il tempo di leggere ciascuna informazione.

\section{Sicurezza e robustezza}
Se un valore sensore e invalido, il display mostra \texttt{ERR} o un messaggio esplicito. Questo evita interpretazioni errate e rende evidenti i guasti.
